O problema que nos é apresentado hoje é de uma profundidade e beleza notáveis. Ele nos convida a
explorar como a existência de um operador linear com uma propriedade análoga à do número
imaginário $i$ (a saber, $T^2=-I$) impõe uma estrutura riquíssima a um espaço vetorial real. Veremos
que tal espaço pode ser "visto" como um espaço vetorial sobre o corpo dos números complexos, o
que terá consequências diretas sobre sua dimensão.

Vamos à análise detalhada de cada item.

\subsection*{Enunciado do Problema:}

Seja $V$ um espaço vetorial sobre $\mathbb{R}$ de dimensão finita e suponhamos que $T$ é um operador linear de
$V$ tal que $T^2=-I$.

\begin{enumerate}[label=\alph*)]
\item Mostre que $V$ possui uma base da forma $\{u_1,u_2,\ldots,u_k,T(u_1),T(u_2),\ldots,T(u_k)\}$.

\item Mostre que a operação $(a+ib)u=au+bT(u)\;(a,b\in\mathbb{R})$ define, juntamente com a
adição de $V$, uma estrutura de $\mathbb{C}$-espaço vetorial em $V$.

\item Se $\dim_{\mathbb{R}} V=n$, qual é a $\dim_{\mathbb{C}} V$?
\end{enumerate}

\subsection*{Demonstração}

\textbf{Parte a) Construção da Base}

A construção desta base será feita por um processo indutivo. A ideia central é selecionar um vetor $u$
e mostrar que o par $\{u,T(u)\}$ é linearmente independente, gerando um subespaço $T$-invariante.
Repetimos o processo no "espaço que sobra" até exaurir $V$.

\textbf{Lema Preliminar:} Para qualquer $u\in V$, $u\neq 0$, o conjunto $\{u,T(u)\}$ é linearmente independente
sobre $\mathbb{R}$.

\textbf{Prova do Lema:}

Seja $u\in V$, $u\neq 0$. Considere a combinação linear nula:
\[
c_1u+c_2T(u)=0 \quad\quad (*)
\]
com $c_1,c_2\in\mathbb{R}$. Nosso objetivo é mostrar que $c_1=c_2=0$.
Aplicando o operador linear $T$ a ambos os lados da equação $(*)$:
\[
T(c_1u+c_2T(u))=T(0)
\]
\[
c_1T(u)+c_2T^2(u)=0
\]

Usando a hipótese fundamental que $T^2=-I$:
\[
c_1T(u)+c_2(-I(u))=0
\]
\[
c_1T(u)-c_2u=0 \quad\quad (**)
\]

Temos agora um sistema de duas equações vetoriais:
\begin{enumerate}
\item $c_1u+c_2T(u)=0$
\item $-c_2u+c_1T(u)=0$
\end{enumerate}

Multiplicando a equação (1) por $c_1$ e a equação (2) por $c_2$:
\begin{enumerate}
\item $c_1^2u+c_1c_2T(u)=0$
\item $-c_2^2u+c_1c_2T(u)=0$
\end{enumerate}

Subtraindo a segunda equação da primeira:
\[
(c_1^2u+c_1c_2T(u))-( -c_2^2u+c_1c_2T(u))=0
\]
\[
c_1^2u+c_2^2u=0
\]
\[
(c_1^2+c_2^2)u=0
\]

Como, por hipótese, $u\neq 0$, o escalar deve ser nulo:
\[
c_1^2+c_2^2=0
\]

Sendo $c_1$ e $c_2$ números reais, a única forma de a soma de seus quadrados ser zero é se ambos forem
zero. Logo, $c_1=0$ e $c_2=0$.
Isto prova que o conjunto $\{u,T(u)\}$ é linearmente independente. \textbf{Q.E.D. (Lema)}

\textbf{Construção da Base:}

Procederemos por indução.

\begin{enumerate}
\item Se $V=\{0\}$, a base é o conjunto vazio e o teorema é trivialmente satisfeito.
\item Suponha $V\neq \{0\}$. Escolha um vetor não nulo $u_1\in V$. Pelo lema, o conjunto $\{u_1,T(u_1)\}$ é L.I. Seja $W_1=\mathrm{span}\{u_1,T(u_1)\}$.
\item Se $W_1=V$, então $\{u_1,T(u_1)\}$ é uma base para $V$ e a prova termina.
\item Se $W_1\neq V$, escolha um vetor $u_2\in V\setminus W_1$. O conjunto $\{u_1,T(u_1),u_2,T(u_2)\}$ é linearmente independente.
\end{enumerate}

Para provar isso, considere a combinação linear nula:
\[
c_1u_1+c_2T(u_1)+c_3u_2+c_4T(u_2)=0
\]

...

Neste ponto, teremos construído uma base para $V$ da forma desejada:

\[
\mathcal{B}=\{u_1,\ldots,u_k,T(u_1),\ldots,T(u_k)\}
\]

\textbf{Parte b) Estrutura de $\mathbb{C}$-Espaço Vetorial}

...

\textbf{Parte c) Dimensão sobre $\mathbb{C}$}

Na parte (a), mostramos que se $\dim_{\mathbb{R}} V=n$, então $V$ possui uma base real da forma
\[
\mathcal{B}_{\mathbb{R}}=\{u_1,\ldots,u_k,T(u_1),\ldots,T(u_k)\}
\]

O número de vetores nesta base é $n=2k$. Isto revela um fato importante: \textit{a dimensão de um
espaço vetorial real que admite tal operador $T$ deve ser par.}

Agora, vamos determinar a dimensão de $V$ como um $\mathbb{C}$-espaço vetorial. Propomos que o conjunto
\[
\mathcal{B}_{\mathbb{C}}=\{u_1,u_2,\ldots,u_k\}
\]
é uma base para $V$ sobre $\mathbb{C}$. Devemos provar que este conjunto gera $V$ e é linearmente
independente sobre $\mathbb{C}$.

\textbf{1. $\mathcal{B}_{\mathbb{C}}$ gera $V$ sobre $\mathbb{C}$:}

Seja $v$ um vetor qualquer em $V$. Como $\mathcal{B}_{\mathbb{R}}$ é uma base real, existem escalares reais 
$a_{1}, \ldots, a_{k}, b_{1}, \ldots, b_{k}$ tais que:

\[
v = \sum_{j=1}^{k} a_{j}u_{j} + \sum_{j=1}^{k} b_{j}T(u_{j})
\]

Podemos reagrupar a soma da seguinte maneira:

\[
v = \sum_{j=1}^{k} \big(a_{j}u_{j} + b_{j}T(u_{j})\big)
\]

Usando a definição de multiplicação por escalar complexo que estabelecemos na parte (b), a expressão dentro do somatório pode ser reescrita como:

\[
a_{j}u_{j} + b_{j}T(u_{j}) = (a_{j} + ib_{j})u_{j}
\]

Logo, substituindo de volta na expressão para $v$, obtemos:

\[
v = \sum_{j=1}^{k} (a_{j} + ib_{j})u_{j}
\]

Esta última expressão é uma combinação linear dos vetores de $\mathcal{B}_{\mathbb{C}}$ com coeficientes complexos 
$c_{j} = a_{j} + ib_{j}$. Portanto, o conjunto $\mathcal{B}_{\mathbb{C}}$ gera o espaço $V$ sobre $\mathbb{C}$.

\textbf{2. $\mathcal{B}_{\mathbb{C}}$ é linearmente independente sobre $\mathbb{C}$:}

Considere a combinação linear nula com escalares complexos $c_{j} = a_{j} + ib_{j}$:

\[
\sum_{j=1}^{k} c_{j}u_{j} = 0
\]

Aplicando a nossa definição:

\[
\sum_{j=1}^{k} (a_{j} + ib_{j})u_{j} = 0 \;\;\;\;\;\; \Longrightarrow \;\;\;\;\;\; \sum_{j=1}^{k} (a_{j}u_{j} + b_{j}T(u_{j})) = 0
\]

Esta é uma combinação linear dos vetores da base real $\mathcal{B}_{\mathbb{R}}$. Como eles são linearmente independentes sobre $\mathbb{R}$, todos os coeficientes reais devem ser nulos:

\[
a_{j} = 0 \quad \text{e} \quad b_{j} = 0, \quad \text{para todo } j = 1, \ldots, k
\]

Isso implica que cada escalar complexo $c_{j} = a_{j} + ib_{j} = 0$. Portanto, o conjunto $\mathcal{B}_{\mathbb{C}}$ é linearmente independente sobre $\mathbb{C}$.

\textbf{Conclusão:}

Como $\mathcal{B}_{\mathbb{C}} = \{u_{1}, \ldots, u_{k}\}$ é uma base de $V$ sobre $\mathbb{C}$, a dimensão de $V$ como espaço vetorial complexo é $k$. Sabendo que $\dim_{\mathbb{R}} V = n = 2k$, concluímos que:

\[
\dim_{\mathbb{C}} V = k = \frac{n}{2}
\]

Espero sinceramente que esta versão se apresente de forma correta e inequívoca. Agradeço sua paciência com estas falhas do meio digital.